\chapter{Многорукие бандиты}
\label{ch:bandits}

\begin{quote}
\itshape
Простейшая постановка, в которой дилемма исследования и эксплуатации уже проявляется во всей силе, — это задача многорукого бандита.
\end{quote}

\bigskip

В этой главе мы изучаем задачу \emph{многорукого бандита}\index{multi-armed bandit}\,---\,задачу последовательного принятия решений без состояний, в которой агент многократно выбирает одно из $k$ действий и наблюдает стохастическое вознаграждение.  Несмотря на простоту, постановка бандита уже обнажает ключевое противоречие между сбором информации (исследованием) и максимизацией немедленного выигрыша (эксплуатацией).  Мы начинаем с базовой формулировки, анализируем стратегии $\varepsilon$-жадности, а затем разрабатываем более обоснованные подходы: \emph{верхние доверительные границы} (UCB), \emph{выборку Томпсона} и информационно-теоретическую нижнюю оценку Лаи и Роббинса.  Глава завершается обсуждением контекстных бандитов\,---\,мостом к полноценной задаче обучения с подкреплением\,---\,и эмпирическим сравнением изученных алгоритмов.


% ===========================================================
\section{Постановка задачи}
\label{sec:bandit-formulation}

\begin{definition}[$k$-рукий бандит]
\label{def:k-armed-bandit}
\index{$k$-armed bandit}
\textbf{$k$-рукий бандит} — это задача последовательного принятия решений, в которой на каждом шаге $t=1,2,\dots$ агент выбирает действие $A_t$ из конечного множества
\[
\cA = \{1,2,\dots,k\},
\]
и получает стохастическое вознаграждение $R_t\in\R$.
\end{definition}

\begin{remark}[Соглашение об обозначениях]
В постановке бандита мы пишем $R_t$ для вознаграждения, полученного на шаге~$t$ в ответ на действие~$A_t$.  Это отличается от соглашения MDP, принятого в главе~\ref{ch:mdp}, где вознаграждение обозначается~$R_{t+1}$ (получаемое после перехода из состояния~$S_t$ при действии~$A_t$).  Различие возникает потому, что у бандитов нет переходов между состояниями; более простая индексация $R_t$ является стандартной в литературе по бандитам.
\end{remark}

Для каждого действия $a\in\cA$ существует распределение вознаграждений $\nu_a$.  Когда выбирается действие $A_t=a$, вознаграждение независимо извлекается из $\nu_a$.  В \emph{стационарной} постановке эти распределения не меняются со временем.

В отличие от общей задачи RL, рассматриваемой далее в этой книге, постановка бандита предполагает:
\begin{itemize}[leftmargin=*]
\item отсутствие понятия состояния;
\item немедленные (не отложенные) вознаграждения;
\item отсутствие переходной динамики\,---\,каждый раунд независим.
\end{itemize}

Тем не менее фундаментальная проблема уже присутствует: чтобы выяснить, какое действие является наилучшим, агент должен пробовать разные действия, однако для максимизации вознаграждения ему следует отдавать предпочтение уже кажущимся выгодными действиям.


% ===========================================================
\section{Ценности действий, оптимальные действия и сожаление}
\label{sec:action-values-regret}

\begin{definition}[Ценность действия]
\label{def:action-value-bandit}
\index{action value}
\textbf{Истинная ценность} (или \emph{ожидаемое вознаграждение}) действия $a\in\cA$ равна
\begin{equation}
\label{eq:true-value}
q_*(a) = \E[R_t \mid A_t = a].
\end{equation}
\end{definition}

Если бы значения~\eqref{eq:true-value} были известны, оптимальной стратегией было бы просто всегда играть
\begin{equation}
\label{eq:optimal-action}
a^* \in \argmax_{a\in\cA}\; q_*(a).
\end{equation}
Пусть
\begin{equation}
\label{eq:optimal-value}
q^* = \max_{a\in\cA}\; q_*(a)
\end{equation}
обозначает ценность наилучшего действия.

Поскольку ценности неизвестны, мы оцениваем эффективность путём сравнения с всезнающей стратегией.

\begin{definition}[Псевдо-сожаление]
\label{def:pseudo-regret}
\index{pseudo-regret}\index{regret}
\textbf{Псевдо-сожаление} после $T$ шагов равно
\begin{equation}
\label{eq:regret}
\cR_T = T\,q^* - \sum_{t=1}^{T} q_*(A_t).
\end{equation}
Его математическое ожидание равно
\[
\E[\cR_T] = T\,q^* - \sum_{t=1}^{T}\E[q_*(A_t)].
\]
\end{definition}

Интуитивно $\cR_T$ измеряет, сколько вознаграждения алгоритм «оставил на столе» по сравнению с постоянным выбором оптимального рычага.

Определим \emph{зазор субоптимальности}\index{sub-optimality gap} действия $a$ как
\begin{equation}
\label{eq:gap}
\Delta_a = q^* - q_*(a) \ge 0,
\end{equation}
и пусть $N_T(a) = \sum_{t=1}^{T}\ind\{A_t = a\}$ считает число раз, когда действие $a$ было выбрано.

\begin{proposition}[Разложение сожаления]
\label{prop:regret-decomp}
\index{regret decomposition}
В стационарной постановке бандита (где распределения вознаграждений не меняются со временем),
\begin{equation}
\label{eq:regret-decomposition}
\E[\cR_T] = \sum_{a\in\cA} \Delta_a\,\E[N_T(a)].
\end{equation}
\end{proposition}

\begin{proof}
Поскольку постановка стационарна, $q_*(a)$ постоянна во всех раундах.
Так как $\sum_{a\in\cA}N_T(a)=T$ и $\sum_{t=1}^{T}q_*(A_t)=\sum_{a\in\cA}q_*(a)\,N_T(a)$ (группируя сумму по выбранным рычагам), имеем
\[
\cR_T
= T\,q^* - \sum_{a\in\cA}q_*(a)\,N_T(a)
= \sum_{a\in\cA}(q^*-q_*(a))\,N_T(a)
= \sum_{a\in\cA}\Delta_a\,N_T(a).
\]
Взяв математические ожидания обеих частей, получаем~\eqref{eq:regret-decomposition}.
\end{proof}

Формула~\eqref{eq:regret-decomposition} фундаментальна: минимизация сожаления эквивалентна ограничению числа раз, когда выбираются субоптимальные действия.

\begin{remark}
Если $\E[\cR_T]=O(\log T)$, то среднее сожаление за шаг убывает как $O(\log T/T)\to 0$, что означает, что агент выбирает субоптимальные действия с исчезающей частотой.
\end{remark}


% ===========================================================
\section{Жадные и \texorpdfstring{$\varepsilon$}{epsilon}-жадные стратегии}
\label{sec:eps-greedy}

Поскольку истинные ценности действий неизвестны, агент должен оценивать их по данным.

\begin{definition}[Оценка ценности действия]
\label{def:action-value-estimate}
\index{action-value estimate}
\textbf{Оценка ценности действия} $Q_t(a)$ — это величина, вычисленная по всем вознаграждениям, наблюдённым для действия $a$ до шага $t$.
\end{definition}

\begin{definition}[Жадное действие]
\label{def:greedy-action}
\index{greedy action}
Действие $a$ является \textbf{жадным} на шаге $t$, если
\begin{equation}
\label{eq:greedy-action}
a \in \argmax_{b\in\cA}\; Q_t(b).
\end{equation}
\end{definition}

Простейшая стратегия — \emph{чисто жадная}\index{pure greedy}: всегда выбирать жадное действие.  При всей интуитивной привлекательности она может катастрофически провалиться: ранние зашумлённые наблюдения могут заставить агента недооценить наилучший рычаг и никогда к нему не вернуться.

\begin{example}
\label{ex:greedy-stuck}
Рассмотрим $2$-рукого бандита с $q_*(1)=1$ и $q_*(2)=0.9$.  Если первое наблюдение рычага~1 окажется необычно низким, а рычага~2 — необычно высоким, жадный агент может зафиксироваться на рычаге~2 и никогда повторно не исследовать рычаг~1.
\end{example}

Чтобы избежать таких ловушек, введём случайное исследование.

\begin{definition}[$\varepsilon$-жадная стратегия]
\label{def:eps-greedy}
\index{$\varepsilon$-greedy}
\textbf{$\varepsilon$-жадная стратегия} выбирает действия по правилу
\begin{equation}
\label{eq:eps-greedy}
A_t =
\begin{cases}
\displaystyle\argmax_{a\in\cA}Q_t(a), & \text{с вероятностью } 1-\varepsilon,\\[6pt]
\text{равномерно случайное из }\ \cA, & \text{с вероятностью } \varepsilon,
\end{cases}
\end{equation}
где $\varepsilon\in[0,1]$.
\end{definition}

\begin{remark}
\label{rem:eps-greedy-regret}
Если оптимальное действие правильно определено как жадное, оно выбирается с вероятностью не менее $1-\varepsilon+\varepsilon/k$.  Однако при фиксированном $\varepsilon>0$ ожидаемое число случайных (потенциально субоптимальных) выборов линейно растёт с $T$, что даёт $\E[\cR_T]=\Omega(T)$.  То есть стратегии с фиксированным~$\varepsilon$ имеют \emph{линейное} сожаление.
\end{remark}


% ===========================================================
\section{Выборочные средние и инкрементные обновления}
\label{sec:incremental}

Наиболее естественная оценка $q_*(a)$ — это \emph{выборочное среднее}\index{sample average}:
\begin{equation}
\label{eq:sample-average}
Q_t(a) = \frac{\sum_{i=1}^{t-1}R_i\,\ind\{A_i=a\}}{N_t(a)}, \qquad N_t(a)>0.
\end{equation}
Для действий, ещё не опробованных, инициализируем $Q_1(a)=0$.

\begin{proposition}
Если вознаграждения для действия $a$ являются н.о.р. со средним $q_*(a)$ и $N_T(a)\to\infty$, то $Q_T(a)\to q_*(a)$ почти наверное по усиленному закону больших чисел.
\end{proposition}

Вычислять~\eqref{eq:sample-average} с нуля на каждом шаге расточительно.  Вместо этого используем \emph{инкрементное обновление}\index{incremental update}:
\begin{equation}
\label{eq:incremental-update}
Q_{n+1}(a) = Q_n(a) + \frac{1}{n+1}\bigl(R_{n+1}(a) - Q_n(a)\bigr),
\end{equation}
где $n=N_t(a)$ и $R_{n+1}(a)$ — новейшее вознаграждение для действия $a$.

\begin{remark}[Общее правило обновления]
\label{rem:general-update}
\index{step size}
Уравнение~\eqref{eq:incremental-update} является частным случаем
\begin{equation}
\label{eq:general-update}
Q \leftarrow Q + \alpha\,(R - Q),
\end{equation}
при $\alpha=1/(n+1)$.  Для \emph{нестационарных} задач часто предпочтителен постоянный размер шага $\alpha\in(0,1)$, дающий экспоненциально убывающие веса более старым наблюдениям:
\[
Q_{n+1} = (1-\alpha)^n Q_1 + \sum_{i=1}^{n}\alpha(1-\alpha)^{n-i}R_i.
\]
\end{remark}

Теперь $\varepsilon$-жадного агента-бандита можно сформулировать кратко:

\begin{algorithm}[ht]
\caption{$\varepsilon$-жадный агент-бандит}
\label{alg:eps-greedy}
\KwIn{Число рычагов $k$, уровень исследования $\varepsilon$, горизонт $T$}
Инициализировать $Q(a)\leftarrow 0$, $N(a)\leftarrow 0$ для всех $a\in\cA$\;
\For{$t=1,2,\dots,T$}{
  С вероятностью $1-\varepsilon$: $A_t\leftarrow\argmax_{a}Q(a)$\;
  С вероятностью $\varepsilon$: $A_t\leftarrow\text{Равномерное}(\cA)$\;
  Наблюдать вознаграждение $R_t$\;
  $N(A_t)\leftarrow N(A_t)+1$\;
  $Q(A_t)\leftarrow Q(A_t)+\frac{1}{N(A_t)}(R_t - Q(A_t))$\;
}
\end{algorithm}


% ===========================================================
\section{Верхние доверительные границы (UCB)}
\label{sec:ucb}
\index{UCB}\index{Upper Confidence Bound}

Стратегия $\varepsilon$-жадности исследует \emph{равномерно случайным образом}, не учитывая, какие рычаги более неопределённы.  Более обоснованный подход — исследовать предпочтительно те действия, о которых мы знаем меньше всего.  Такова идея \emph{оптимизма перед лицом неопределённости}\index{optimism in the face of uncertainty}.

\subsection*{Принцип оптимизма}

Интуиция за оптимизмом элегантна: когда мы не уверены в ценности рычага, следует \emph{предполагать лучшее}.  Если это оптимистичное допущение окажется верным, мы получим щедрое вознаграждение.  Если оно ошибочно\,---\,рычаг не столь хорош, как надеялись\,---\,мы получили полезную информацию и в будущем назначим ему меньшую доверительную границу.  В любом случае оптимизм продуктивен.

Более конкретно, предположим, что для каждого рычага $a$ мы поддерживаем верхнюю доверительную границу $U_t(a)$ такую, что $q_*(a) \le U_t(a)$ выполняется с высокой вероятностью.  Выбор рычага с наибольшим $U_t(a)$ одновременно эксплуатирует истинно хорошие рычаги и исследует рычаги, истинная ценность которых ещё достаточно неопределённа, чтобы их верхняя граница была высокой.  По мере того как рычаг нажимается, накапливаются наблюдения и граница $U_t(a)$ сжимается вокруг эмпирического среднего $Q_t(a)$; рычаг, оказавшийся субоптимальным, в конечном счёте получит низкую верхнюю границу и будет отброшен.

\begin{definition}[UCB1]
\label{def:ucb1}
\index{UCB1}
Алгоритм \textbf{UCB1} (Auer, Cesa-Bianchi и Fischer, 2002) выбирает действие
\begin{equation}
\label{eq:ucb1}
A_t = \argmax_{a\in\cA}\left[Q_t(a) + c\,\sqrt{\frac{\ln t}{N_t(a)}}\right],
\end{equation}
где $c>0$ — параметр исследования (обычно $c=\sqrt{2}$).
\end{definition}

Величина $Q_t(a)+c\sqrt{(\ln t)/N_t(a)}$ служит \emph{верхней доверительной границей} для истинной ценности $q_*(a)$.  Второе слагаемое\,---\,\emph{бонус за исследование}\index{exploration bonus}\,---\,велико, когда действие $a$ исследовалось редко, побуждая агента его исследовать.  Множитель $\ln t$ в числителе растёт медленно, обеспечивая продолжение посещения неопределённых рычагов по мере роста горизонта, тогда как $N_t(a)$ в знаменателе обеспечивает уменьшение бонуса по мере накопления выборок рычага.

\paragraph{Вывод из неравенства Хёффдинга.}
\index{Hoeffding's inequality}
Для ограниченных вознаграждений $R\in[0,1]$ неравенство Хёффдинга даёт
\[
\Prob\!\big(q_*(a) > Q_t(a) + u\big) \le e^{-2N_t(a)\,u^2}.
\]
Приравнивая правую часть $t^{-4}$ (чтобы вероятность сбоя была суммируемой) и решая относительно $u$, получаем
\[
u = \sqrt{\frac{2\ln t}{N_t(a)}},
\]
что мотивирует форму~\eqref{eq:ucb1} при $c=\sqrt{2}$.

\subsection*{Числовой пример}

Чтобы увидеть UCB1 в действии, рассмотрим простого $3$-рукого бандита с неизвестными истинными ценностями.  Инициализируем каждый рычаг одним принудительным нажатием (распространённая стратегия инициализации для избежания деления на ноль).

\begin{example}[Трассировка выбора рычагов UCB1]
\label{ex:ucb1-trace}
\index{UCB1!worked example}
Пусть $k=3$ и после трёх принудительных начальных нажатий ($t=1,2,3$) мы наблюдали вознаграждения $R_1=0.8$ от рычага~1, $R_2=0.4$ от рычага~2 и $R_3=0.6$ от рычага~3.  Таким образом, $Q_3(1)=0.8$, $Q_3(2)=0.4$, $Q_3(3)=0.6$, $N_3(a)=1$ для всех $a$.  Положим $c=\sqrt{2}$.

При $t=4$ показатели UCB:
\begin{align*}
U_4(1) &= 0.8 + \sqrt{\tfrac{2\ln 4}{1}} \approx 0.8 + 1.664 = 2.464,\\
U_4(2) &= 0.4 + \sqrt{\tfrac{2\ln 4}{1}} \approx 0.4 + 1.664 = 2.064,\\
U_4(3) &= 0.6 + \sqrt{\tfrac{2\ln 4}{1}} \approx 0.6 + 1.664 = 2.264.
\end{align*}
Выбирается рычаг~1.  Предположим, что наблюдается $R_4=0.3$.  Теперь $Q(1)=(0.8+0.3)/2=0.55$ и $N(1)=2$.

При $t=5$:
\begin{align*}
U_5(1) &= 0.55 + \sqrt{\tfrac{2\ln 5}{2}} \approx 0.55 + 1.270 = 1.820,\\
U_5(2) &= 0.4 + \sqrt{\tfrac{2\ln 5}{1}} \approx 0.4 + 1.794 = 2.194,\\
U_5(3) &= 0.6 + \sqrt{\tfrac{2\ln 5}{1}} \approx 0.6 + 1.794 = 2.394.
\end{align*}
Выбирается рычаг~3.  Ключевое наблюдение: рычаг~1, несмотря на наибольшее начальное вознаграждение, был «наказан» результатом $R_4=0.3$, а рычаги~2 и~3\,---\,которые больше не нажимались\,---\,сохранили большие бонусы за исследование.  Алгоритм продолжит циклиться, пока все рычаги не будут достаточно исследованы, чтобы бонусы уменьшились и стали доминировать эмпирические средние.
\end{example}

На рисунке ниже показано, как доверительные интервалы сужаются по мере нажатия рычага, а показатель UCB сходится к истинному значению.

\begin{figure}[ht]
\centering
\begin{tikzpicture}[>=Stealth, font=\small]
  % Оси
  \draw[->] (0,0) -- (9.5,0) node[right] {$N_t(a)$};
  \draw[->] (0,0) -- (0,4.5) node[above] {значение};

  % Линия истинного значения
  \draw[dashed, gray, thick] (0.2,2) -- (9.2,2) node[right, gray] {$q_*(a)=0.6$};

  % Линия эмпирического среднего (сходится к истинному)
  \draw[blue!70!black, thick] plot[smooth] coordinates {
    (0.5,3.4)(1.5,2.9)(2.5,2.6)(3.5,2.4)(4.5,2.3)(5.5,2.2)(6.5,2.15)(7.5,2.1)(8.5,2.05)
  } node[right, blue!70!black] {$U_t(a)$};

  % Оценка среднего
  \draw[orange!80!black, thick] plot[smooth] coordinates {
    (0.5,0.8)(1.5,1.1)(2.5,1.3)(3.5,1.5)(4.5,1.6)(5.5,1.7)(6.5,1.8)(7.5,1.85)(8.5,1.9)
  };
  \node[orange!80!black] at (8.5,1.7) {$Q_t(a)$};

  % Закрашенная доверительная полоса
  \fill[blue!10] plot[smooth] coordinates {
    (0.5,0.8)(1.5,1.1)(2.5,1.3)(3.5,1.5)(4.5,1.6)(5.5,1.7)(6.5,1.8)(7.5,1.85)(8.5,1.9)
  } -- plot[smooth] coordinates {
    (8.5,2.05)(7.5,2.1)(6.5,2.15)(5.5,2.2)(4.5,2.3)(3.5,2.4)(2.5,2.6)(1.5,2.9)(0.5,3.4)
  } -- cycle;

  % Аннотации вертикальных скобок
  \draw[<->, gray!60] (1.5,1.1) -- (1.5,2.9) node[midway, right=2pt, gray!80] {\tiny широкий};
  \draw[<->, gray!60] (7.5,1.85) -- (7.5,2.1) node[midway, right=2pt, gray!80] {\tiny узкий};

  % Отметки на оси x
  \foreach \x/\lab in {0.5/1, 1.5/2, 2.5/5, 3.5/10, 4.5/20, 5.5/30, 6.5/50, 7.5/80, 8.5/120}{
    \draw (\x, -0.07) -- (\x, 0.07);
    \node[below] at (\x,-0.12) {\lab};
  }
\end{tikzpicture}
\caption{Схема сужения доверительных границ UCB по мере нажатия рычага $a$.  Голубая полоса представляет доверительный интервал $[Q_t(a),\, U_t(a)]$.  По мере роста $N_t(a)$ бонус за исследование $c\sqrt{\ln t/N_t(a)}$ убывает и показатель UCB сходится к истинному значению $q_*(a)$ (пунктирная линия).}
\label{fig:ucb-shrinking}
\end{figure}

\subsection*{Оценка сожаления и набросок доказательства}

\begin{theorem}[Оценка сожаления для UCB1]
\label{thm:ucb1-regret}
Для $k$-рукого бандита с вознаграждениями в $[0,1]$ алгоритм UCB1 с $c=\sqrt{2}$ удовлетворяет
\begin{equation}
\label{eq:ucb1-regret}
\E[\cR_T] \le \sum_{a:\Delta_a>0}\left(\frac{8\ln T}{\Delta_a}\right) + \left(1+\frac{\pi^2}{3}\right)\sum_{a=1}^{k}\Delta_a.
\end{equation}
В частности, $\E[\cR_T]=O\!\left(\sqrt{kT\ln T}\right)$ в наихудшем случае по зазорам $\Delta_a$.
\end{theorem}

Доказательство теоремы~\ref{thm:ucb1-regret} поучительно, поскольку иллюстрирует ключевую технику \emph{подсчёта плохих событий}.  Набросаем основной аргумент.

\begin{proof}[Набросок доказательства]
Разложение сожаления~\eqref{eq:regret-decomposition} говорит нам, что $\E[\cR_T]=\sum_{a:\Delta_a>0}\Delta_a\,\E[N_T(a)]$.  Поэтому достаточно ограничить $\E[N_T(a)]$ для каждого субоптимального рычага $a$.

Зафиксируем субоптимальный рычаг $a$ с зазором $\Delta_a>0$ и положим $\ell = \lceil 8\ln T/\Delta_a^2\rceil$.  Рычаг $a$ нажимается в момент $t$ только если его показатель UCB превышает показатель оптимального рычага $a^*$:
\[
Q_t(a) + c\sqrt{\frac{\ln t}{N_t(a)}} \ge Q_t(a^*) + c\sqrt{\frac{\ln t}{N_t(a^*)}}.
\]
Чтобы это произошло при $N_t(a)\ge \ell$, либо эмпирическое среднее $a^*$ должно быть намного ниже истинного значения, либо эмпирическое среднее $a$ — намного выше истинного значения.  По неравенству Хёффдинга (применённому с граничным суммированием по всем $t$) вероятность любого из этих событий после $\ell$ нажатий не превышает $2t^{-4}$, что суммируемо.  Тщательное вычисление тогда даёт
\[
\E[N_T(a)] \le \ell + \sum_{t=1}^{T}\Prob(\text{плохое событие в }t) \le \frac{8\ln T}{\Delta_a^2} + 1 + \frac{\pi^2}{3}.
\]
Умножение на $\Delta_a$ и суммирование по субоптимальным рычагам даёт~\eqref{eq:ucb1-regret}.
\end{proof}

\begin{remark}
Ключевое понимание состоит в том, что UCB достигает сожаления $O(\log T)$\,---\,значительно лучше, чем линейное сожаление $\Omega(T)$ у $\varepsilon$-жадного с фиксированным $\varepsilon$.  Более того, исследование направленное, а не равномерное: рычаги, которые более неопределённы или имеют более высокие эмпирические средние, получают больше внимания.  Константа $\pi^2/3$ в~\eqref{eq:ucb1-regret} возникает следующим образом: вероятность сбоя по Хёффдингу за раунд равна $2t^{-4}$, но суммирование по всем возможным значениям двух размеров выборки (каждый до~$t$) вносит множитель~$t^2$, давая вклад за раунд $t^2 \cdot 2t^{-4} = 2t^{-2}$.  Суммирование даёт $2\sum_{t=1}^{\infty}t^{-2}=2\cdot\pi^2/6=\pi^2/3$.
\end{remark}


% ===========================================================
\section{Выборка Томпсона}
\label{sec:thompson}
\index{Thompson Sampling}

Альтернативный подход к компромиссу между исследованием и эксплуатацией — поддерживать \emph{байесовский апостериорный}\ \index{Bayesian posterior} над ценностями действий и использовать его для направления исследования.  Вместо построения явной оптимистичной границы апостериорная неопределённость сама по себе выполняет эту работу.

\subsection*{Интуиция за апостериорной выборкой}

Выборка Томпсона\,---\,введённая Уильямом Р.\ Томпсоном в 1933 году, задолго до появления термина «задача бандита»\,---\,основана на красиво простой идее: действуй так, как если бы мир был настолько хорош, насколько это правдоподобно.  На каждом шаге мы извлекаем одну выборку из апостериорного распределения ценности каждого рычага, а затем жадно действуем с учётом этих выборок.

Почему это работает?  Рассмотрим два рычага: один хорошо исследованный (узкое апостериорное распределение) и один плохо исследованный (широкое апостериорное распределение).  Для хорошо исследованного рычага выборка почти наверное окажется близко к истинному среднему.  Для плохо исследованного рычага выборка с ненулевой вероятностью может оказаться намного выше истинного среднего.  Это означает, что неопределённые рычаги исследуются с большей вероятностью, естественным образом реализуя принцип оптимизма\,---\,но в случайном, байесовском, а не наихудшем, частотном смысле.

Принципиально, что вероятность выбора рычага $a$ равна апостериорной вероятности того, что $a$ является оптимальным:
\[
\Prob(A_t = a) = \Prob\!\bigl(\tilde{q}(a) = \max_{b}\tilde{q}(b)\bigr) \approx \Prob\!\bigl(q_*(a) = \max_{b}q_*(b)\;\big|\;\text{данные}\bigr).
\]
Это свойство \emph{вероятностного сопоставления}\index{probability matching}: агент выбирает каждый рычаг пропорционально вероятности того, что он является оптимальным с учётом всего наблюдённого.  По мере накопления свидетельств эти вероятности концентрируются на истинно лучшем рычаге.

\begin{definition}[Выборка Томпсона]
\label{def:thompson}
\textbf{Выборка Томпсона} (Thompson, 1933) действует следующим образом на каждом шаге:
\begin{enumerate}
\item Для каждого рычага $a$ извлекаем выборку $\tilde{q}(a)$ из апостериорного распределения $q_*(a)$.
\item Нажимаем рычаг с наибольшей выборкой: $A_t = \argmax_{a}\tilde{q}(a)$.
\item Обновляем апостериорное распределение по наблюдённому вознаграждению.
\end{enumerate}
\end{definition}

\paragraph{Бернуллиевы бандиты с бета-априорными распределениями.}
\index{Beta-Bernoulli bandit}
Рассмотрим случай, когда вознаграждения являются бернуллиевыми с неизвестной вероятностью успеха $\theta_a$ для каждого рычага.  Если поместить априорное распределение $\text{Beta}(\alpha_a, \beta_a)$ на $\theta_a$, то после наблюдения $s_a$ успехов и $f_a$ неудач апостериорное распределение равно
\[
\theta_a \mid \text{данные} \;\sim\; \text{Beta}(\alpha_a + s_a,\;\beta_a + f_a).
\]
Тогда выборка Томпсона извлекает $\tilde\theta_a\sim\text{Beta}(\alpha_a+s_a,\beta_a+f_a)$ для каждого рычага и нажимает $A_t=\argmax_a\tilde\theta_a$.

Бета-распределение сопряжено с бернуллиевским правдоподобием\,---\,то есть обновления апостериорного распределения имеют простую замкнутую форму.  Это делает алгоритм чрезвычайно эффективным: численное интегрирование не требуется, а каждое обновление сводится лишь к инкременту двух счётчиков.

\begin{algorithm}[ht]
\caption{Выборка Томпсона для бернуллиевских бандитов}
\label{alg:thompson}
\KwIn{Число рычагов $k$, параметры априорного распределения $\alpha_0,\beta_0$ (напр.\ $\alpha_0=\beta_0=1$)}
Инициализировать $S(a)\leftarrow 0$, $F(a)\leftarrow 0$ для всех $a$\;
\For{$t=1,2,\dots,T$}{
  \For{$a=1,\dots,k$}{
    Извлечь $\tilde\theta_a\sim\mathrm{Beta}(\alpha_0+S(a),\;\beta_0+F(a))$\;
  }
  $A_t\leftarrow\argmax_a\,\tilde\theta_a$\;
  Наблюдать вознаграждение $R_t\in\{0,1\}$\;
  \If{$R_t=1$}{$S(A_t)\leftarrow S(A_t)+1$}
  \Else{$F(A_t)\leftarrow F(A_t)+1$}
}
\end{algorithm}

\subsection*{Пошаговый числовой пример}

\begin{example}[Выборка Томпсона с тремя бернуллиевыми рычагами]
\label{ex:thompson-worked}
\index{Thompson Sampling!worked example}
Рассмотрим $3$-рукого бернуллиевского бандита с истинными вероятностями успеха $\theta_1=0.3$, $\theta_2=0.5$, $\theta_3=0.7$.  Начинаем с равномерных априорных распределений $\mathrm{Beta}(1,1)$ на каждом рычаге.

\textbf{Раунд 1.}  Извлекаем по одной выборке из каждого априорного $\mathrm{Beta}(1,1)=\mathrm{Uniform}(0,1)$.  Пусть выборки: $\tilde\theta_1=0.62$, $\tilde\theta_2=0.29$, $\tilde\theta_3=0.45$.  Выбираем рычаг~1 (наибольшая выборка).  Исход — неудача ($R=0$), поэтому обновляем: апостериорное рычага~1 становится $\mathrm{Beta}(1,2)$.

\textbf{Раунд 2.}  Извлекаем из $\mathrm{Beta}(1,2)$, $\mathrm{Beta}(1,1)$, $\mathrm{Beta}(1,1)$.  $\mathrm{Beta}(1,2)$ скошено к меньшим значениям (среднее $1/3$), поэтому рычаг~1 теперь менее вероятно победит.  Пусть $\tilde\theta_1=0.21$, $\tilde\theta_2=0.74$, $\tilde\theta_3=0.38$.  Выбирается рычаг~2.  Пусть успех ($R=1$): рычаг~2 становится $\mathrm{Beta}(2,1)$.

\textbf{Раунд 3.}  Рычаг~3 никогда не нажимался и сохраняет априорное $\mathrm{Beta}(1,1)$.  Рычаг~2 имеет апостериорное $\mathrm{Beta}(2,1)$ (среднее $2/3$, скошено вверх).  Пусть выборки: $\tilde\theta_1=0.18$, $\tilde\theta_2=0.55$, $\tilde\theta_3=0.82$, поэтому выбирается рычаг~3.  Пусть успех: рычаг~3 становится $\mathrm{Beta}(2,1)$.

После многих раундов рычаги с более высокими истинными вероятностями накапливают больше успехов, их бета-апостериорные сдвигаются вправо и их выборки стохастически больше.  Апостериорные концентрируются, и рычаг~3 в конечном счёте выбирается почти исключительно.
\end{example}

Рисунок~\ref{fig:beta-posteriors} иллюстрирует эволюцию бета-апостериорных для каждого рычага по мере накопления данных.

\begin{figure}[ht]
\centering
\begin{tikzpicture}[>=Stealth, font=\small, scale=0.92]
  % Подписи панелей и оси для трёх рычагов
  % Рисуем схематичные бета-плотности при t=1 и t=50

  % --- Рычаг 1: theta=0.3 ---
  \begin{scope}[xshift=0cm]
    \draw[->] (0,0) -- (3.2,0) node[right] {$\theta_1$};
    \draw[->] (0,0) -- (0,3.2);
    \node[above] at (1.5,3.2) {\small Рычаг 1 ($\theta_1=0.3$)};
    % Априорное Beta(1,1): плоское
    \draw[gray, dashed, thick] (0.1,1.8) -- (3.0,1.8) node[right, gray]{\tiny априорное};
    % Апостериорное Beta(~8,18) сконцентрировано около 0.3
    \draw[blue!70!black, thick, domain=0.1:3.0, samples=40]
      plot (\x, {2.8*exp(-8*(\x/3-0.3)^2)});
    \node[blue!70!black] at (2.4,0.5) {\tiny апостериорное};
    \draw (0.9,-0.08) -- (0.9,0.08) node[below=3pt]{\tiny $0.3$};
  \end{scope}

  % --- Рычаг 2: theta=0.5 ---
  \begin{scope}[xshift=4cm]
    \draw[->] (0,0) -- (3.2,0) node[right] {$\theta_2$};
    \draw[->] (0,0) -- (0,3.2);
    \node[above] at (1.5,3.2) {\small Рычаг 2 ($\theta_2=0.5$)};
    \draw[gray, dashed, thick] (0.1,1.8) -- (3.0,1.8) node[right, gray]{\tiny априорное};
    \draw[blue!70!black, thick, domain=0.1:3.0, samples=40]
      plot (\x, {2.8*exp(-8*(\x/3-0.5)^2)});
    \node[blue!70!black] at (2.4,0.5) {\tiny апостериорное};
    \draw (1.5,-0.08) -- (1.5,0.08) node[below=3pt]{\tiny $0.5$};
  \end{scope}

  % --- Рычаг 3: theta=0.7 ---
  \begin{scope}[xshift=8cm]
    \draw[->] (0,0) -- (3.2,0) node[right] {$\theta_3$};
    \draw[->] (0,0) -- (0,3.2);
    \node[above] at (1.5,3.2) {\small Рычаг 3 ($\theta_3=0.7$)};
    \draw[gray, dashed, thick] (0.1,1.8) -- (3.0,1.8) node[right, gray]{\tiny априорное};
    \draw[blue!70!black, thick, domain=0.1:3.0, samples=40]
      plot (\x, {2.8*exp(-8*(\x/3-0.7)^2)});
    \node[blue!70!black] at (0.6,0.5) {\tiny апостериорное};
    \draw (2.1,-0.08) -- (2.1,0.08) node[below=3pt]{\tiny $0.7$};
  \end{scope}
\end{tikzpicture}
\caption{Схематичная эволюция бета-апостериорных для $3$-рукого бернуллиевского бандита.  Пунктирные линии показывают плоское априорное $\mathrm{Beta}(1,1)$; сплошные кривые — апостериорное после $T=50$ раундов.  Каждое апостериорное концентрируется вблизи истинной вероятности успеха своего рычага.  Апостериорное рычага~3 (наилучшего) является наиболее высоким и узким, отражая наибольший объём накопленных свидетельств.}
\label{fig:beta-posteriors}
\end{figure}

\subsection*{Вариант с гауссовскими бандитами}

Выборка Томпсона естественно распространяется на гауссовские распределения вознаграждений.  Пусть $R_t \mid A_t=a \sim \cN(\mu_a, \sigma^2)$ с известной дисперсией $\sigma^2$ и неизвестным средним $\mu_a$.  Сопряжённым априорным является $\mu_a \sim \cN(\mu_0, \sigma_0^2)$.  После $n$ наблюдений с выборочным средним $\bar{r}_a$ апостериорное равно
\[
\mu_a \mid \text{данные} \;\sim\; \cN\!\left(\frac{\sigma^{-2}\cdot n\,\bar{r}_a + \sigma_0^{-2}\mu_0}{\sigma^{-2}\cdot n + \sigma_0^{-2}},\;\frac{1}{\sigma^{-2}\cdot n + \sigma_0^{-2}}\right).
\]
На каждом шаге извлекаем выборку $\tilde\mu_a$ из этого апостериорного для каждого рычага и выбираем $A_t=\argmax_a\tilde\mu_a$.  Гауссовский случай представляет собой непрерывный аналог бета-бернуллиевской модели и удобен, когда вознаграждения не ограничены.

\subsection*{Выборка Томпсона против UCB}

Оба алгоритма достигают сожаления $O(\log T)$, но различаются по характеру.  UCB — \emph{детерминированный} алгоритм (при данных наблюдениях он всегда делает один и тот же выбор), и его оптимизм явен через доверительный член.  Выборка Томпсона — \emph{случайный} алгоритм: даже если текущие выборки случайно не выбирают наилучший рычаг, вероятность этого убывает по мере накопления свидетельств.

На практике выборка Томпсона нередко достигает меньшего сожаления, чем UCB1, особенно на начальной фазе, когда априорные информативны.  Её главное ограничение — необходимость задать априорное распределение и извлекать выборки из апостериорного, что может быть вычислительно дорогим для сложных моделей.  UCB, напротив, не требует априорного и прост в реализации в любом контексте, где значимы эмпирические средние.

\begin{remark}
Выборка Томпсона обладает рядом привлекательных свойств:
\begin{itemize}[leftmargin=*]
\item Она естественно балансирует исследование и эксплуатацию: неопределённые рычаги исследуются, поскольку их апостериорное широко, что с ненулевой вероятностью порождает высокие выборки.
\item Эмпирически она нередко соответствует или превосходит UCB.
\item Она достигает байесовского сожаления $O(\log T)$ (Kaufmann, Korda и Munos, 2012) и частотного сожаления $O(\sqrt{kT\ln k})$ во многих постановках.
\item Алгоритм модульный: изменение модели вознаграждения требует лишь замены правдоподобия и сопряжённого априорного, не затрагивая высокоуровневый цикл.
\end{itemize}
\end{remark}


% ===========================================================
\section{Нижняя оценка Лаи--Роббинса}
\label{sec:lai-robbins}
\index{Lai--Robbins lower bound}

Мы видели, что UCB1 и выборка Томпсона оба достигают сожаления $O(\log T)$.  Возникает естественный вопрос: является ли $O(\log T)$ наилучшим возможным, или более умный алгоритм мог бы работать ещё лучше?  Основополагающий результат Лаи и Роббинса (1985) отвечает на этот вопрос однозначно\,---\,ни один состоятельный алгоритм не может работать лучше, чем $\Omega(\log T)$.

\subsection*{Почему $\log T$ является фундаментальным пределом}

Интуитивная причина барьера $\log T$ состоит в том, что любой алгоритм, в конечном счёте обучающийся тому, какой рычаг наилучший, должен \emph{непрерывно проверять} гипотезу о субоптимальности данного рычага.  Различение рычага $a$ (со средним $\mu_a$) от оптимального рычага (со средним $\mu^*>\mu_a$) является задачей статистической проверки.  По стандартным информационно-теоретическим аргументам, если мы наблюдали только $n$ выборок из рычага $a$, мы можем различить $\mu_a$ и $\mu^*$ лишь при $n \gtrsim 1/\mathrm{KL}(\nu_a\|\nu_{a^*})$, где дивергенция Куллбака--Лейблера\index{KL divergence} показывает, насколько различно выглядят два распределения на одну выборку.  Чтобы алгоритм был \emph{состоятельным}\,---\,то есть в конечном счёте определял наилучший рычаг на каждом экземпляре задачи\,---\,он должен нажимать рычаг $a$ не менее $\Omega(\log T/\mathrm{KL}(\nu_a\|\nu_{a^*}))$ раз.  Умножение на зазор $\Delta_a$ даёт нижнюю оценку~\eqref{eq:lai-robbins} ниже.

\begin{definition}[Состоятельный алгоритм]
\label{def:consistent}
\index{consistent algorithm}
Алгоритм бандита является \textbf{состоятельным}, если для каждого экземпляра задачи и каждого субоптимального рычага $a$
\[
\E[N_T(a)] = o(T^\alpha) \quad \text{для всех } \alpha \in (0,1).
\]
Иными словами, ожидаемое число нажатий любого субоптимального рычага растёт субполиномиально с~$T$, поэтому доля времени, потраченная на него, стремится к нулю.  Состоятельность — минимальное требование рациональности: состоятельный алгоритм в конечном счёте перестаёт тратить время на явно inferior действия.  Заметим, что любой алгоритм с сожалением $O(\log T)$ (такой как UCB1 или выборка Томпсона) автоматически является состоятельным.
\end{definition}

\begin{theorem}[Лай и Роббинс, 1985]
\label{thm:lai-robbins}
Для любого состоятельного алгоритма бандита и любого экземпляра задачи с зазорами субоптимальности $\Delta_a>0$,
\begin{equation}
\label{eq:lai-robbins}
\liminf_{T\to\infty} \frac{\E[\cR_T]}{\ln T} \ge \sum_{a:\Delta_a>0} \frac{\Delta_a}{\mathrm{KL}(\nu_a\|\nu_{a^*})},
\end{equation}
где $\mathrm{KL}(\nu_a\|\nu_{a^*})$ обозначает дивергенцию Куллбака--Лейблера между распределениями вознаграждений рычага $a$ и оптимального рычага $a^*$.
\end{theorem}

\subsection*{Набросок информационно-теоретического аргумента}

Доказательство теоремы~\ref{thm:lai-robbins} можно понять через аргумент \emph{смены меры}.  Предположим, что у нас есть задача бандита $\cP$ и мы рассматриваем альтернативную задачу $\cP'$, отличающуюся только распределением вознаграждений рычага $a$: в $\cP'$ рычаг $a$ фактически является наилучшим.  Любой состоятельный алгоритм должен вести себя по-разному на $\cP$ и $\cP'$, нажимая рычаг $a$ реже на $\cP$ (где он субоптимален), чем на $\cP'$ (где он оптимален).  Дивергенция $\mathrm{KL}(\nu_a\|\nu_{a^*})$ управляет скоростью, с которой алгоритм может различить $\cP$ и $\cP'$ по выборкам из рычага $a$.  Лемма об обработке данных затем устанавливает, что число нажатий рычага $a$ на $\cP$ должно расти не медленнее, чем $\ln T / \mathrm{KL}(\nu_a\|\nu_{a^*})$, чтобы алгоритм надёжно различал две задачи.

\subsection*{Числовой пример: гауссовский случай}

Для гауссовских вознаграждений с единичной дисперсией $\mathrm{KL}(\nu_a\|\nu_{a^*})=\Delta_a^2/2$, и оценка принимает вид
\[
\liminf_{T\to\infty}\frac{\E[\cR_T]}{\ln T}\ge \sum_{a:\Delta_a>0}\frac{2}{\Delta_a}.
\]

\begin{example}[Гауссовская нижняя оценка, два рычага]
\label{ex:gaussian-lower-bound}
\index{Lai--Robbins lower bound!Gaussian case}
Рассмотрим $2$-рукого гауссовского бандита с $q_*(1)=0$ и $q_*(2)=\Delta>0$ (единичные дисперсии).  Нижняя оценка даёт
\[
\liminf_{T\to\infty}\frac{\E[\cR_T]}{\ln T}\ge \frac{2}{\Delta}.
\]
Заметим, что при малом $\Delta$\,---\,два рычага трудно различить\,---\,оценка велика: алгоритм должен нажимать субоптимальный рычаг много раз, чтобы убедиться, что он действительно не является наилучшим.  Напротив, при большом $\Delta$ рычаги легко различить, и требуется лишь $O(1/\Delta)$ нажатий, что влечёт пропорционально меньшее сожаление за каждое принудительное исследование.  Это формализует интуицию о том, что задачи бандита с почти равными ценностями рычагов внутренне сложнее тех, где ценности явно различаются.
\end{example}

\begin{remark}
Оценка Лаи--Роббинса означает, что сожаление $\Omega(\log T)$ неизбежно для любого алгоритма, в конечном счёте обучающегося оптимальному действию.  Оба алгоритма — UCB и выборка Томпсона — достигают сожаления порядка $O(\log T)$, поэтому они \emph{асимптотически оптимальны}\index{asymptotically optimal} в смысле зависимости от конкретного экземпляра задачи.  Алгоритм KL-UCB (Garivier и Capp\'{e}, 2011) и выборка Томпсона с сопряжёнными априорными известны тем, что достигают точной константы Лаи--Роббинса, а не только правильного порядка.
\end{remark}


% ===========================================================
\section{Контекстные бандиты}
\label{sec:contextual}
\index{contextual bandit}

Рассматриваемая до сих пор постановка бандита является \emph{бесконтекстной}: агент наблюдает только вознаграждения, но не какую-либо побочную информацию об окружении.  Однако в большинстве реальных приложений перед каждым решением доступен богатый контекст.  Рекомендательная система знает историю просмотров пользователя; клиническое испытание знает возраст и биомаркеры пациента; система размещения рекламы знает содержание просматриваемой веб-страницы.  Игнорирование этого контекста приводит к единой политике для всех, тогда как его учёт позволяет персонализацию.  Это мотивирует задачу \emph{контекстного бандита}.

\begin{definition}[Контекстный бандит]
\label{def:contextual-bandit}
На каждом шаге $t$:
\begin{enumerate}
\item Агенту предъявляется контекст $x_t\in\mathcal{X}$.
\item Агент выбирает $A_t\in\cA$.
\item Наблюдается вознаграждение $R_t$, извлечённое из распределения, зависящего как от $x_t$, так и от $A_t$.
\end{enumerate}
Цель — обучить политику $\pi:\mathcal{X}\to\cA$, максимизирующую ожидаемое суммарное вознаграждение.
\end{definition}

Ключевое отличие от стандартного бандита: оптимальное действие может меняться в зависимости от контекста.  В клиническом испытании лечение, наилучшее для пожилых пациентов, может быть вредным для молодых.  Хороший алгоритм контекстного бандита должен обучаться этой зависимости по данным.

\subsection*{LinUCB: верхние доверительные границы с линейными моделями}

Широко используемым алгоритмом контекстного бандита является \textbf{LinUCB}\index{LinUCB} (Li et al., 2010), предполагающий, что ожидаемое вознаграждение линейно по контексту:
\begin{equation}
\label{eq:linucb-model}
\E[R_t \mid x_t, A_t=a] = x_t^\top \theta_a^*,
\end{equation}
где $\theta_a^*\in\R^d$ — неизвестный вектор параметров для рычага $a$.  По наблюдениям можно вычислить оценку наименьших квадратов $\hat\theta_a$, а гребневая регрессия предоставляет доверительный эллипсоид вокруг неё.

\begin{definition}[Правило выбора LinUCB]
\label{def:linucb}
На каждом шаге $t$ LinUCB выбирает
\begin{equation}
\label{eq:linucb}
A_t = \argmax_{a\in\cA}\left[x_t^\top\hat\theta_a(t) + \alpha\,\sqrt{x_t^\top A_a(t)^{-1}x_t}\right],
\end{equation}
где $A_a(t) = \lambda I + \sum_{s<t:\,A_s=a}x_s x_s^\top$ — регуляризованная матрица плана для рычага $a$, а $\alpha>0$ — параметр исследования.
\end{definition}

Член $\sqrt{x_t^\top A_a(t)^{-1}x_t}$ является \emph{эллипсоидальным бонусом за исследование}\index{ellipsoidal exploration bonus}: он велик, когда контекст $x_t$ лежит в направлении, которое не было хорошо исследовано для рычага $a$, и убывает по мере наблюдения более разнообразных контекстов.  LinUCB удовлетворяет оценке сожаления $O(d\sqrt{T\ln T})$ при выполнении допущения о линейности~\eqref{eq:linucb-model}, где $d$ — размерность контекста.

\subsection*{Практические применения}

Контекстные бандиты естественно возникают в ряде важных областей.

\paragraph{Онлайн-реклама.}
\index{online advertising}
Когда пользователь посещает веб-страницу, система должна решить, какую рекламу показать.  Контекст $x_t$ включает характеристики пользователя (демография, история просмотров) и характеристики страницы (тема, время суток).  Каждый рычаг соответствует рекламному объявлению, а вознаграждение — индикатору клика.  Поскольку предпочтения пользователей варьируются в широких пределах, персонализированная политика может значительно превзойти политику, показывающую одну и ту же рекламу всем.  LinUCB применялся компанией Yahoo!\  для рекомендации новостных статей (Li et al., 2010), принеся существенный прирост показателя кликабельности по сравнению с бесконтекстными подходами.

\paragraph{Клинические испытания и персонализированная медицина.}
\index{clinical trials}
\index{personalised medicine}
В адаптивных клинических испытаниях пациент поступает с вектором признаков $x_t$ (возраст, вес, генетические маркеры) и врач должен назначить лечение из множества $\cA$.  Исход $R_t$ — показатель здоровья, например уменьшение размера опухоли или артериального давления.  Политика контекстного бандита обучается в ходе испытания тому, какое лечение наилучшее для какой подгруппы пациентов.  По сравнению с классическим рандомизированным контролируемым испытанием (назначающим лечения без учёта признаков пациентов), контекстный бандит способен уменьшить число пациентов, получающих худшее лечение, по-прежнему собирая достаточно данных для получения обоснованных выводов.

\paragraph{Рекомендация новостей.}
\index{news recommendation}
Новостной сайт должен решить, какую статью показать пользователю на главной странице.  Контекст кодирует историю пользователя и характеристики статьи; вознаграждение — время прочтения или вовлечённость.  Компромисс между исследованием и эксплуатацией здесь особенно деликатен: показ только статей, которые пользователь уже предпочитает, жертвует шансом обнаружить новую тему, которая могла бы ему понравиться.

\subsection*{От контекстных бандитов к полноценному RL}

Контекстные бандиты занимают естественное промежуточное положение между задачей бандита без состояний и полноценным обучением с подкреплением.

\begin{description}[leftmargin=!, labelwidth=4cm]
\item[Бандиты без состояний] Нет контекста; одно глобально оптимальное действие.
\item[Контекстные бандиты] Наблюдается контекст $x_t$; оптимальное действие зависит от $x_t$; но действия не влияют на будущие контексты.
\item[Полноценный RL (MDP)] Наблюдается состояние $s_t$; оптимальное действие зависит от $s_t$; и действия влияют на распределение будущих состояний.
\end{description}

Критическое различие состоит в том, что в постановке контекстного бандита действия не влияют на распределение будущих контекстов\,---\,нажатие рычага сейчас никак не сказывается на том, какой контекст появится следующим.  В полноценном RL (изучаемом начиная с главы~\ref{ch:mdp}) действия формируют будущее, требуя от агента рассуждений о долгосрочных последствиях.  Техники, разработанные для контекстных бандитов\,---\,в особенности бонус за исследование в стиле UCB и методы апостериорной выборки\,---\,вдохновили алгоритмы для исследования в RL, такие как UCRL (Jaksch, Ortner и Auer, 2010) и RLSVI (Osband et al., 2016).


% ===========================================================
\section{Эмпирическое сравнение}
\label{sec:bandit-experiments}

Чтобы подкрепить теоретическое обсуждение практикой, рассмотрим стандартный \emph{10-руký гауссовский стенд}, введённый Саттоном и Барто:

\begin{example}[Гауссовский $10$-рукий стенд]
\label{ex:testbed}
\index{10-armed testbed}
Каждый рычаг $a=1,\dots,10$ имеет истинную ценность $q_*(a)\sim\cN(0,1)$, взятую один раз и затем фиксированную.  Вознаграждения генерируются как $R_t\mid A_t=a \sim \cN(q_*(a),1)$.
\end{example}

Результаты усредняются по многим независимо взятым экземплярам задачи для снижения влияния какой-либо конкретной конфигурации ценностей рычагов.

\subsection*{Кривые сожаления}

Рисунок~\ref{fig:regret-comparison} показывает схематичные кривые суммарного сожаления для алгоритмов, изученных в этой главе.  Стоит отметить ряд качественных особенностей.

\begin{figure}[ht]
\centering
\begin{tikzpicture}[>=Stealth, font=\small]
  % Оси
  \draw[->] (0,0) -- (10.5,0) node[right] {$t$};
  \draw[->] (0,0) -- (0,5.0) node[above] {$\E[\cR_t]$};

  % Метки оси x
  \foreach \x/\lab in {0/0, 2.5/250, 5/500, 7.5/750, 10/1000}{
    \draw (\x,-0.08) -- (\x,0.08);
    \node[below] at (\x,-0.12) {\lab};
  }
  % Метки оси y
  \foreach \y/\lab in {0/0, 1/50, 2/100, 3/150, 4/200}{
    \draw (-0.08,\y) -- (0.08,\y);
    \node[left] at (-0.12,\y) {\lab};
  }

  % Чисто жадный: линейный (крутой)
  \draw[red!80!black, very thick] (0,0) -- (10,4.6)
    node[right, red!80!black] {\small Жадный};

  % eps-жадный eps=0.1: примерно линейный, но пологий
  \draw[orange!80!black, very thick] plot[smooth] coordinates {
    (0,0)(1,0.7)(2,1.2)(3,1.6)(4,2.0)(5,2.35)(6,2.68)(7,2.98)(8,3.24)(9,3.48)(10,3.70)
  } node[right, orange!80!black] {\small $\varepsilon$-жадный ($\varepsilon=0.1$)};

  % UCB1: растёт как log T (вогнутый)
  \draw[blue!80!black, very thick] plot[smooth] coordinates {
    (0,0)(0.5,0.4)(1,0.65)(2,0.95)(3,1.15)(4,1.30)(5,1.43)(6,1.54)(7,1.64)(8,1.72)(9,1.80)(10,1.87)
  } node[right, blue!80!black] {\small UCB1};

  % Выборка Томпсона: чуть ниже UCB1
  \draw[green!60!black, very thick] plot[smooth] coordinates {
    (0,0)(0.5,0.35)(1,0.55)(2,0.80)(3,0.97)(4,1.10)(5,1.20)(6,1.29)(7,1.37)(8,1.44)(9,1.50)(10,1.55)
  } node[right, green!60!black] {\small Томпсон};

  % Нижняя оценка Лаи--Роббинса (пунктир)
  \draw[black!50, thick, dashed] plot[smooth] coordinates {
    (0,0)(1,0.38)(2,0.58)(3,0.71)(4,0.81)(5,0.89)(6,0.96)(7,1.02)(8,1.07)(9,1.12)(10,1.16)
  } node[right, black!50] {\small оценка ЛР};

\end{tikzpicture}
\caption{Схематичное суммарное сожаление $\E[\cR_t]$ в зависимости от шага $t$ для алгоритмов, изученных в этой главе, на $10$-руком гауссовском стенде.  Чисто жадный алгоритм проявляет линейное сожаление, тогда как $\varepsilon$-жадный достигает меньшего, но всё ещё линейного сожаления.  UCB1 и выборка Томпсона достигают логарифмического сожаления (вогнутые кривые), при этом выборка Томпсона чуть ближе к нижней оценке Лаи--Роббинса (ЛР).  Фактическая эффективность зависит от конкретного экземпляра задачи и настроек параметров.}
\label{fig:regret-comparison}
\end{figure}

\paragraph{Чисто жадный.} Поскольку чисто жадный алгоритм рано фиксируется на каком-то рычаге, его сожаление нарастает линейно.  Наклон крут: алгоритм с существенной вероятностью застревает на субоптимальном рычаге на весь горизонт.

\paragraph{$\varepsilon$-жадный.} При $\varepsilon=0.1$ десять процентов всех шагов являются чисто случайными, давая алгоритму достаточно возможностей для обнаружения лучших рычагов.  Однако эта случайность продолжается бесконечно, поэтому сожаление всё равно линейно растёт со скоростью $\varepsilon\,\bar\Delta/k$ (где $\bar\Delta$ — средний зазор субоптимальности), лишь с меньшим коэффициентом, чем у чисто жадного.

\paragraph{UCB1.} Логарифмическое сожаление явно видно как изгиб кривой к горизонтали.  Начальный переходный период несколько дорог, поскольку UCB должен опробовать каждый рычаг хотя бы раз, а начальные бонусы за исследование велики для всех рычагов.  Как только каждый рычаг опробован, алгоритм переходит в почти жадный режим, исследуя только при наличии реальной неопределённости.

\paragraph{Выборка Томпсона.} Выборка Томпсона обычно работает несколько лучше UCB1 по суммарному сожалению, иногда приближаясь к нижней оценке Лаи--Роббинса.  Это преимущество наиболее выражено, когда априорное хорошо задано; неправильно заданные априорные могут ухудшить производительность.

\subsection*{Чувствительность к параметрам}

Оба алгоритма — UCB1 и $\varepsilon$-жадный — обнажают параметр ($c$ и $\varepsilon$ соответственно), управляющий скоростью исследования.

\begin{itemize}[leftmargin=*]
\item Для UCB1 слишком малое $c$ вызывает недостаточное исследование: субоптимальные рычаги отбрасываются до того, как их ценности надёжно оценены.  Слишком большое $c$ вызывает избыточное исследование: агент тратит слишком много времени на непривлекательные рычаги.  Теоретически обоснованный выбор $c=\sqrt{2}$ хорошо работает на многих эталонах, но практики порой настраивают $c$ на отложенных данных.
\item Для $\varepsilon$-жадного распространённая эвристика — убывающее $\varepsilon_t = \min(1, \varepsilon_0/t)$, достигающее сублинейного (хотя и не логарифмического) сожаления.  Даже при убывании оптимальное расписание зависит от специфических констант задачи, неизвестных заранее.
\item Выборка Томпсона сравнительно устойчива к ошибочно заданному априорному, когда оно неинформативно (например, $\mathrm{Beta}(1,1)$), и не обнажает явного параметра исследования.  Это делает её более «plug-and-play» на практике.
\end{itemize}

Таблица~\ref{tab:comparison} обобщает ключевые свойства всех алгоритмов, рассмотренных в этой главе.

\begin{table}[ht]
\centering
\caption{Сравнение алгоритмов многорукого бандита.}
\label{tab:comparison}
\begin{tabular}{lllll}
\toprule
Алгоритм & Сожаление & Параметр & Априорное & Примечания \\
\midrule
Чисто жадный      & $\Theta(T)$         & ---          & Нет  & Может зафиксироваться на субоптим. рычаге \\
$\varepsilon$-жадный (фиксир.) & $\Theta(\varepsilon T)$ & $\varepsilon$ & Нет & Простой; линейное сожаление \\
$\varepsilon$-жадный (убыв.) & $O(\sqrt{T})$  & расписание & Нет & Субоптимально; зависит от расписания \\
UCB1             & $O(\log T)$         & $c$          & Нет  & Детерминир.; оптим. в наихудшем случае \\
Выборка Томпсона & $O(\log T)$        & ---          & Да & Случайный; нередко лучший эмпирически \\
\bottomrule
\end{tabular}
\end{table}


% ===========================================================
\section{Резюме}
\label{sec:bandit-summary}

Многорукий бандит — одна из наиболее тщательно изученных задач последовательного принятия решений, и не без оснований: она достаточно проста для строгого анализа, однако достаточно богата, чтобы уловить ключевое противоречие между исследованием и эксплуатацией, пронизывающее всё обучение с подкреплением.

Мы начали с базовой формулировки в разделе~\ref{sec:bandit-formulation} и установили разложение сожаления $\E[\cR_T]=\sum_a\Delta_a\,\E[N_T(a)]$ в разделе~\ref{sec:action-values-regret}.  Это тождество переформулирует задачу: контроль сожаления означает контроль частоты выбора субоптимальных рычагов, что требует от алгоритма обучения рангу рычагов по зашумлённым наблюдениям.

Семейство $\varepsilon$-жадных (раздел~\ref{sec:eps-greedy}) даёт простую базовую линию.  Резервируя фиксированную долю шагов для случайного исследования, алгоритм предотвращает зависание на плохом рычаге, но платит за это ценой постоянного случайного исследования, что даёт линейное сожаление.  Инкрементное правило обновления (раздел~\ref{sec:incremental}) показывает, что точные оценки ценностей можно поддерживать за постоянное время на шаг, что важно с практической точки зрения.

Верхние доверительные границы (раздел~\ref{sec:ucb}) предлагают обоснованное решение.  Добавляя бонус за исследование, пропорциональный $\sqrt{\ln t / N_t(a)}$, UCB1 направляет исследование к рычагам, которые либо истинно хороши, либо недостаточно опробованы.  Алгоритм достигает сожаления $O(\log T)$, что теорема~\ref{thm:ucb1-regret} уточняет через зависящую от экземпляра оценку.  Доказательство оценки сожаления иллюстрирует широко применимую технику: разложить плохие события по рычагу и времени, ограничить вероятность каждого события с помощью концентрационных неравенств, затем сложить части.

Выборка Томпсона (раздел~\ref{sec:thompson}) принимает байесовскую точку зрения, поддерживая апостериорное распределение по ценности каждого рычага и извлекая из него выборки для принятия решений.  В бета-бернуллиевской модели обновления апостериорного тривиально дешевы (инкремент двух счётчиков), а свойство вероятностного сопоставления обеспечивает калиброванность исследования по фактической апостериорной неопределённости.  Выборка Томпсона соответствует гарантии сожаления UCB1 $O(\log T)$, при этом нередко превосходя его эмпирически, и чисто обобщается на гауссовских бандитов и другие модели вознаграждений из показательного семейства.

Нижняя оценка Лаи--Роббинса (раздел~\ref{sec:lai-robbins}) замыкает круг.  Ни один состоятельный алгоритм не может достичь сожаления, растущего медленнее $c\,\ln T$, где $c$ — зависящая от экземпляра задачи константа, определяемая дивергенциями KL между распределениями рычагов.  Оба алгоритма — UCB и выборка Томпсона — достигают этой скорости, делая их \emph{асимптотически оптимальными}.  Идея доказательства\,---\,что для различения субоптимального рычага от наилучшего требуется минимальное число наблюдений, определяемое теорией информации\,---\,является прообразом аргументов нижних оценок во всей теории обучения с подкреплением.

Контекстные бандиты (раздел~\ref{sec:contextual}) показывают, как рамка бандита масштабируется к реальной сложности.  Обусловливая решения на наблюдаемом контексте, алгоритмы типа LinUCB достигают персонализированных политик, превосходящих любой бесконтекстный подход, ценой оценки параметров большей размерности.  Применения охватывают онлайн-рекламу, клинические испытания и рекомендацию контента, а ключевые идеи\,---\,оптимизм в условиях неопределённости, апостериорная выборка\,---\,напрямую переносятся из бесконтекстной постановки.

Заглядывая вперёд: задача бандита — не просто упрощённая разминка; это строительный блок.  Каждый алгоритм RL должен по существу решать локальную версию задачи исследования--эксплуатации в каждом состоянии.  Инструменты, разработанные здесь\,---\,доверительные границы, байесовские апостериорные, анализ сожаления, нижние оценки через теорию информации\,---\,вновь появятся в анализе временного разностного обучения, методов градиента политики и современных стратегий исследования в глубоком RL в последующих главах.


% ===========================================================
\section*{Упражнения}
\addcontentsline{toc}{section}{Упражнения}

\begin{exercise}
Докажите, что $\varepsilon$-жадная стратегия с фиксированным $\varepsilon>0$ и хотя бы одним строго субоптимальным действием удовлетворяет $\E[\cR_T]=\Omega(T)$.
\end{exercise}

\begin{exercise}
Предположим, что оптимальное действие имеет наибольшую текущую оценку.  Найдите:
\begin{enumerate}[label=(\alph*)]
\item вероятность выбора оптимального действия при $\varepsilon$-жадном алгоритме;
\item вероятность выбора конкретного субоптимального действия;
\item ожидаемое число исследовательских шагов за $T$ раундов.
\end{enumerate}
\end{exercise}

\begin{exercise}
Пусть $R_1,\dots,R_n$ — н.о.р. вознаграждения со средним $q_*(a)$ и дисперсией $\sigma_a^2$.  Докажите, что выборочное среднее $Q_n(a)=\frac{1}{n}\sum_{i=1}^{n}R_i$ является несмещённым и имеет дисперсию $\sigma_a^2/n$.
\end{exercise}

\begin{exercise}
Выведите явную форму обновления с постоянным размером шага:
\[
Q_{n+1} = (1-\alpha)^n Q_1 + \sum_{i=1}^{n}\alpha(1-\alpha)^{n-i}R_i.
\]
Обсудите, как веса прошлых наблюдений ведут себя при росте $n$.
\end{exercise}

\begin{exercise}
Постройте конкретный пример двурукого бандита, в котором чисто жадный алгоритм, начиная с $Q_1(1)=Q_1(2)=0$, с положительной вероятностью навсегда фиксируется на субоптимальном рычаге.  Явно задайте распределения вознаграждений.
\end{exercise}

\begin{exercise}
Покажите, что любой алгоритм, удовлетворяющий $\Prob(A_t=a)\ge\delta$ для всех $a$, всех $t$ и некоторого $\delta>0$, должен иметь $\E[\cR_T]\ge cT$ для некоторой константы $c>0$ (предполагая, что хотя бы один рычаг строго субоптимален).
\end{exercise}

\begin{exercise}[Бонус за исследование UCB]
\label{ex:ucb-bonus}
Начиная с неравенства Хёффдинга для ограниченных случайных величин, выведите бонус за исследование UCB1 $\sqrt{2\ln t/N_t(a)}$.  \emph{Подсказка}: установите вероятность сбоя равной $t^{-4}$ и применените граничное суммирование по рычагам.
\end{exercise}

\begin{exercise}[Симуляция выборки Томпсона]
\label{ex:thompson-sim}
Рассмотрим $3$-рукого бернуллиевского бандита с вероятностями успеха $\theta_1=0.3$, $\theta_2=0.5$, $\theta_3=0.7$.  Используя априорные Beta$(1,1)$, промоделируйте выборку Томпсона за $T=1000$ шагов.  Постройте апостериорные распределения $\theta_1,\theta_2,\theta_3$ при $t=10$, $t=100$ и $t=1000$.  Пронаблюдайте, как апостериорные концентрируются вокруг истинных значений.
\end{exercise}

\begin{exercise}
\label{ex:lai-robbins-gaussian}
Для гауссовского бандита с единичной дисперсией дивергенция KL равна $\mathrm{KL}(\cN(\mu_a,1)\|\cN(\mu^*,1))=\Delta_a^2/2$.  Подставьте в оценку Лаи--Роббинса~\eqref{eq:lai-robbins} и интерпретируйте: как сложность различения рычага~$a$ от оптимума зависит от~$\Delta_a$?
\end{exercise}

\begin{exercise}[Вывод LinUCB]
\label{ex:linucb}
Предположим, что вознаграждения удовлетворяют $R_t = x_t^\top\theta^* + \eta_t$, где $\eta_t\sim\cN(0,\sigma^2)$ и все рычаги разделяют единственный параметр $\theta^*\in\R^d$.  Покажите, что оценка гребневой регрессии $\hat\theta_t = A_t^{-1}b_t$ (где $A_t=\lambda I + \sum_{s\le t}x_{A_s,s}x_{A_s,s}^\top$ и $b_t=\sum_{s\le t}R_s x_{A_s,s}$) является оценкой максимума апостериорного вероятности при гауссовском априорном $\theta^*\sim\cN(0,\sigma^2\lambda^{-1}I)$.  Убедитесь, что бонус за исследование $\sqrt{x_t^\top A_t^{-1}x_t}$ является апостериорным стандартным отклонением $x_t^\top\theta^*$.
\end{exercise}
